\section{Introduction}
\label{sec:introduction}
\label{sec:intro}

Speech-driven facial animation aims to generate temporally synchronized facial motion from speech audio, supporting applications in digital avatars and embodied interaction.
Recent advances in portrait animation~\cite{ji2021audio,zhang2023sadtalker}, 3D facial animation~\cite{fan2022faceformer,xing2023codetalker}, and audio-conditioned neural rendering~\cite{guo2021ad} have improved the quality of generated talking faces.
In parallel, compact facial representations, including 3D morphable models, FLAME, and blendshape parameterizations, provide practical control spaces for reconstruction and animation~\cite{blanz2023morphable,li2017learning,menzel2022automated}.
Within blendshape-based representations, ARKit defines semantically meaningful facial action coefficients~\cite{apple_arkit}, offering an interpretable interface between speech articulation and mouth motion.

Existing audio-driven methods predict facial motion from learned acoustic representations, leaving linguistic and articulatory structure implicit in the input waveform~\cite{cudeiro2019voca,richard2021meshtalk,fan2022faceformer,park2023said}.
Viseme-based approaches explicitly model phonetic groups~\cite{zhou2018visemenet}, and more recent methods incorporate textual or linguistic cues~\cite{fan2022joint,xu2024kmtalk}.
However, we observe that relatively little research has explored how aligned audio, transcripts, and phonemes can jointly guide a multimodal language model to generate interpretable ARKit motion trajectories.
This leads us to ask: \emph{Can a multimodal language model use explicit linguistic and articulatory structure to generate accurate mouth-motion trajectories from speech?}

Bearing this in mind, we consider the connection between phonetic units and the facial actions involved in their production.
Bilabial consonants involve lip closure, open vowels generally require greater jaw opening, and transitions between phonemes shape the evolution of mouth motion.
These articulatory relations suggest a natural connection between linguistic inputs and the action semantics of ARKit coefficients.
Recent language-aware facial models provide further motivation: SemanticFace uses language-aligned semantic supervision to estimate facial actions from images~\cite{kang2026semanticface}, while KeyframeFace incorporates language-model priors and semantic keyframes for text-driven facial animation~\cite{wu2025keyframeface}.
Together, these connections motivate us to make transcripts, phonemes, and articulation guidance explicit inputs to a model that generates facial controls. 
Realizing this idea requires both local alignment and suitable training objectives.
Long recordings must be organized into units in which linguistic content and motion targets share temporal boundaries.
Furthermore, low-valued coefficients dominate the training targets and can bias token-level learning.
Even with suitable conditioning, next-token supervision does not explicitly evaluate the temporal dynamics or geometric effects of the resulting trajectories.
These considerations motivate a framework that combines aligned inputs, balanced coefficient generation, and motion-level feedback.

To this end, we present AudioMouth, a language-assisted framework that formulates the prediction of 32 mouth-related ARKit blendshape trajectories as schema-constrained generation with a multimodal language model.
AudioMouth first segments long utterances into semantically coherent units and uses shared temporal boundaries to extract aligned audio, transcript, phoneme, and coefficient subsequences.
The generator then predicts fixed-schema coefficient sequences under explicit phoneme-to-articulation guidance.
A distribution-balanced supervised objective reweights coefficient-value tokens to reduce the dominance of frequently occurring low-valued targets.
To connect token-level learning with the quality of generated motion, we further refine the generator using Group Relative Policy Optimization (GRPO).
The reward compares sampled trajectories with reference motion in terms of coefficient reconstruction, velocity, acceleration, and mouth geometry.
For geometric feedback, we introduce a simplified lip-geometry discrepancy (SLMD) that maps coefficient differences to weighted mouth-vertex displacements through a fixed blendshape basis.
This provides a 3D geometric surrogate without rendering or landmark detection and complements coefficient-wise and temporal supervision. Our contributions are threefold:
\begin{itemize}
    \item We propose AudioMouth, a language-assisted framework that integrates aligned linguistic inputs and explicit articulation guidance into the structured generation of mouth-related ARKit trajectories with a multimodal language model.

    \item We develop a three-stage framework that combines audio-linguistic unit construction, distribution-balanced supervised generation, and motion-aware policy refinement. We introduce SLMD to provide geometric feedback directly from coefficient-induced mouth deformation.

    \item On the selected data, AudioMouth outperforms four Audio2ARKit baselines across all five reported metrics, reducing Fr\'echet distance from the best baseline value of $25.228$ to $10.820$. Ablations support the benefits of linguistic and articulatory conditioning and show additional improvements in distributional and geometric fidelity after motion-aware refinement.
\end{itemize}
